\chapter{Descriptive set theory: projection theorems and the everywhere-Borel filtration}
\label{ch:dst}

The measurable multiplicative theory of the preceding chapters produces the singular
(``issue \#6'') forward Lyapunov filtration only \emph{almost everywhere}: the a.e.\ route
(Chapter~\ref{chap:lyapunov}) proves the projector \(x\mapsto\mathrm{orthProjMatrix}(V x)\)
\texttt{AEMeasurable} by pushing a measurable graph through the universal measurability of analytic
sets. Upgrading that to \emph{everywhere} Borel measurability --- the payoff of issue \#11 --- is a
genuinely descriptive-set-theoretic problem: the sublevel set \(\{x\mid \mathrm{infDist}\,c\,(V
x)\le r\}\) is the projection to the base of a Borel subset of a product of Polish spaces, and
Borel sets are not closed under projection (their projections are exactly the analytic sets, which
need not be Borel). What rescues the situation is that the fibres one projects are \emph{compact},
and for Borel sets with compact sections the projection \emph{is} Borel --- the \textbf{Novikov
projection theorem}. This chapter formalises the full ladder that reaches it, entirely from
Mathlib's binary Lusin separation, and consumes it for the everywhere-Borel singular filtration.

The ladder is classical (Srivastava, \emph{A Course on Borel Sets}, \S\S4.6--4.7; Kechris,
\emph{Classical Descriptive Set Theory}, \S\S14, 28): the generalized first separation theorem of
Novikov (the \(\N\)-ary generalization of binary Lusin separation) \(\Rightarrow\) the weak
reduction principle for coanalytic sets \(\Rightarrow\) the Saint Raymond and Kunugui--Novikov
section theorems \(\Rightarrow\) the compact-section projection theorem. Notably \emph{no coanalytic
pointclass need be developed}: the whole reduction ladder is run on analytic sets and their Borel
separators. The chapter opens with the classical Lusin theorem in its continuous-on-a-compact form
(a companion descriptive-set-theoretic result, consumed by the measurable Livšic tier), turns the
same analytic / coanalytic vocabulary on a \emph{complexity} question --- the descriptive complexity
of continuous-section existence over a Polish space of continuous-dynamics systems (issue \#61) ---
and closes with an honest statement of the general \(\sigma\)-compact-section theorem
(Arsenin--Kunugui) as a documented research frontier.

Throughout, spaces are Polish (separable completely metrizable) carrying their Borel
\(\sigma\)-algebra, ``analytic'' means a continuous image of Baire space \(\N\to\N\) (Mathlib's
\texttt{MeasureTheory.AnalyticSet}), and everything is formalized sorry-free and audited to the
axiom set \(\{\texttt{propext},\ \texttt{Classical.choice},\ \texttt{Quot.sound}\}\).

\section{Lusin's theorem}

We first record the classical Lusin theorem, a companion descriptive result that Mathlib supplies
the ingredients for but not the assembled statement.

\begin{theorem}[Lusin's theorem, continuous-on-a-compact form]\label{thm:dst-lusin}
  \lean{ErgodicTheory.lusin_continuousOn}\leanok
  Let \(u:X\to\R\) be Borel measurable on a Polish space \(X\) carrying a finite Borel measure
  \(\mu\). For every tolerance \(\varepsilon\neq 0\) there is a \emph{compact} set \(K\subseteq X\)
  with \(\mu(K^{\mathsf c})<\varepsilon\) on which \(u\) is continuous
  (\(\mathrm{ContinuousOn}\,u\,K\)).
\end{theorem}
\begin{proof}\leanok
  The arctan-compression route (Cohn, \emph{Measure Theory}, 2nd ed., Thm.~7.4.4; Rudin,
  \emph{Real and Complex Analysis}, Thm.~2.24), which avoids any integrability hypothesis on
  \(u\). Replace \(u\) by the bounded measurable \(v=\arctan\circ u\), valued in
  \((-\pi/2,\pi/2)\) hence in \(L^1(\mu)\) since \(\mu\) is finite. Bounded continuous functions
  are dense in \(L^1\) (a finite measure on a metrizable Borel space is weakly regular), so pick
  continuous \(g_n\) with \(\mathrm{eLpNorm}(v-g_n)\,1\,\mu\to 0\); \(L^1\)-convergence gives
  convergence in measure, hence an a.e.-convergent subsequence. Egorov's theorem upgrades this to
  \emph{uniform} convergence off a small measurable set, and inner regularity carves a compact
  \(K\) out of its complement with \(\mu(K^{\mathsf c})<\varepsilon\). On \(K\) the continuous
  \(g_n\) converge uniformly to \(v\), so \(v\) is continuous on \(K\); composing with \(\tan\)
  (continuous on \((-\pi/2,\pi/2)\), and \(\tan\circ\arctan=\mathrm{id}\)) recovers
  \(\mathrm{ContinuousOn}\,u\,K\).
\end{proof}

\section{Analytic-set closure lemmas}

The reduction ladder is run on analytic sets, and needs three closure properties that are genuine
gaps in Mathlib's \texttt{AnalyticSet} API. All three are natural upstream candidates for
\texttt{Mathlib.MeasureTheory.Constructions.Polish.Basic} (Srivastava \S4.3; Kechris \S14).

\begin{lemma}[Binary intersection of analytic sets is analytic]\label{thm:dst-inter}
  \lean{MeasureTheory.AnalyticSet.inter'}\leanok
  In a Hausdorff space, if \(s\) and \(t\) are analytic then so is \(s\cap t\).
\end{lemma}
\begin{proof}\leanok
  Write \(s\cap t=\bigcap_{b:\mathrm{Bool}}\mathrm{cond}\,b\,s\,t\) and apply Mathlib's countable
  intersection \texttt{AnalyticSet.iInter} over the two-element index \(\mathrm{Bool}\).
\end{proof}

\begin{lemma}[Binary union of analytic sets is analytic]\label{thm:dst-union}
  \lean{MeasureTheory.AnalyticSet.union'}\leanok
  If \(s\) and \(t\) are analytic then so is \(s\cup t\).
\end{lemma}
\begin{proof}\leanok
  Write \(s\cup t=\bigcup_{b:\mathrm{Bool}}\mathrm{cond}\,b\,s\,t\) and apply the countable union
  \texttt{AnalyticSet.iUnion} over \(\mathrm{Bool}\).
\end{proof}

\begin{lemma}[Product of analytic sets is analytic]\label{thm:dst-prod}
  \lean{MeasureTheory.AnalyticSet.prod'}\leanok
  If \(s\subseteq X\) and \(t\subseteq Y\) are analytic then \(s\times t\subseteq X\times Y\) is
  analytic.
\end{lemma}
\begin{proof}\leanok
  Parametrise \(s=\mathrm{range}\,f\), \(t=\mathrm{range}\,g\) by continuous maps from Polish
  spaces (\texttt{analyticSet\_iff\_exists\_polishSpace\_range}); then \(s\times t\) is the range of
  the continuous product map \(f\times g\) from the product Polish space.
\end{proof}

The a.e.\ singular filtration of Chapter~\ref{chap:lyapunov} is built on a fourth, deeper property
of analytic sets --- their \emph{universal measurability} --- which the everywhere route of this
chapter deliberately bypasses (it produces \emph{null}-measurability, hence only \texttt{AEMeasurable}
projectors). We cross-reference it here for contrast.

\begin{theorem}[Analytic sets are universally measurable]\label{thm:dst-nullmeas}
  \lean{MeasureTheory.AnalyticSet.nullMeasurableSet}\leanok
  Every analytic set in a standard Borel space is \texttt{NullMeasurableSet} for every s-finite
  measure.
\end{theorem}
\begin{proof}\leanok
  Choquet's capacitability theorem (Kechris, Thm.~30.13; Srivastava, Thm.~4.3.1): for a finite
  Borel measure the compact capacity \(\mathrm{compactCap}\,\mu\,s=\sup\{\mu K\mid K\text{ compact},
  K\subseteq s\}\) agrees with \(\mu s\) on analytic sets, so an increasing union of compact
  subsets exhausts \(s\) up to a null set, giving \(s=^{\mu}B\) with \(B\) Borel; the s-finite case
  dominates \(\mu\) by a finite measure. This route needs no analytic complement (Suslin's theorem
  does), which is exactly why it applies to a Borel projection. It underlies the \emph{a.e.}\
  singular tier (\Cref{thm:dst-aemeasurable}); the everywhere upgrade below replaces it by the
  Novikov projection theorem.
\end{proof}

\section{Novikov's generalized first separation theorem}

The engine of the whole development is the \(\N\)-ary generalization, due to Novikov (with the
elegant proof of Mokobodzki), of Mathlib's binary Lusin separation
\texttt{measurablySeparable\_range\_of\_disjoint}. Where the binary theorem separates two disjoint
analytic sets by a Borel set, Novikov's theorem separates a \emph{countable} family with empty
intersection by measurable supersets that still intersect emptily.

\begin{definition}[Borel-separated family]\label{thm:dst-sepfam}
  \lean{ErgodicTheory.SepFam}\leanok
  A countable family \(E:\N\to\mathcal P(X)\) is \emph{Borel separated} (\texttt{SepFam}) if there
  are measurable supersets \(B_n\supseteq E_n\) with empty intersection
  \(\bigcap_n B_n=\emptyset\) (Srivastava's terminology after 4.6.1).
\end{definition}

\begin{lemma}[The single-entry split, Lemma 4.6.2]\label{thm:dst-sepfam-split}
  \lean{ErgodicTheory.sepFam_of_forall_update}\leanok
  \uses{thm:dst-sepfam}
  Fix a position \(j\) and a countable decomposition \(E_j\subseteq\bigcup_m G_m\). If every
  refinement \(\mathrm{update}\,E\,j\,(G_m)\) (replacing the \(j\)-th entry by \(G_m\)) is Borel
  separated, then so is \(E\). Its contrapositive existential form
  \lean{ErgodicTheory.exists_not_sepFam_update}
  states: if \(E\) is \emph{not} separated then some refinement \(\mathrm{update}\,E\,j\,(G_m)\) is
  not separated either.
\end{lemma}
\begin{proof}\leanok
  \uses{thm:dst-sepfam}
  Given separators \(B^{(m)}\) for each refinement, assemble a separator for \(E\): at position
  \(j\) take \(\bigcup_m B^{(m)}_j\), at every other position \(n\) take \(\bigcap_m B^{(m)}_n\).
  These are measurable, contain \(E\) (using \(E_j\subseteq\bigcup_m G_m\) at \(j\)), and have
  empty intersection: a point in the intersection lies in some \(B^{(m_0)}_j\), hence in
  \(\bigcap_n B^{(m_0)}_n=\emptyset\), a contradiction.
\end{proof}

\begin{lemma}[One stage of the Mokobodzki recursion, inductive Lemma 4.6.2]\label{thm:dst-stage-step}
  \lean{ErgodicTheory.stage_step}\leanok
  \uses{thm:dst-sepfam-split}
  Applied to images \(A_n=\mathrm{range}\,f_n\) of continuous maps \(f_n:(\N\to\N)\to X\) from Baire
  space, consider the stage-\(k\) family
  \lean{ErgodicTheory.famE}
  \(\mathrm{famE}\,f\,\sigma\,k=\bigl(f_n[\mathrm{cylinder}(\sigma_n,k-n)]\bigr)_n\) of
  cylinder images around centres \(\sigma_n\). If \(\mathrm{famE}\,f\,\sigma\,k\) is not Borel
  separated, there are refined centres \(\sigma'\), each within the length-\((k-n)\) cylinder of
  \(\sigma_n\), for which \(\mathrm{famE}\,f\,\sigma'\,(k+1)\) is still not separated.
\end{lemma}
\begin{proof}\leanok
  \uses{thm:dst-sepfam-split}
  Extend the entries one coordinate at a time along the mixed-length intermediate family
  \lean{ErgodicTheory.famMix}
  \(\mathrm{famMix}\), by an inner induction over the position \(j\le k+1\): decompose the
  \(j\)-th cylinder into the one-longer cylinders \(\bigcup_m\mathrm{cylinder}(\mathrm{update}\,
  \sigma'_j\,(k-j)\,m)\), and invoke the single-entry split
  (\Cref{thm:dst-sepfam-split}) to pick a coordinate value \(m_0\) preserving non-separatedness.
  At \(j=k+1\) the mixed family is exactly \(\mathrm{famE}\,f\,\sigma'\,(k+1)\).
\end{proof}

\begin{theorem}[First separation for ranges, Srivastava 4.6.1]\label{thm:dst-sepfam-ranges}
  \lean{ErgodicTheory.sepFam_ranges}\leanok
  \uses{thm:dst-stage-step}
  If \(f_n:(\N\to\N)\to X\) are continuous with \(\bigcap_n\mathrm{range}\,f_n=\emptyset\), then the
  family \((\mathrm{range}\,f_n)\) is Borel separated.
\end{theorem}
\begin{proof}\leanok
  \uses{thm:dst-stage-step, thm:dst-sepfam}
  Suppose not. Iterating one stage of the recursion (\Cref{thm:dst-stage-step}) from the trivial
  centres builds, by dependent choice, centres refining coordinate by coordinate; a stationarity
  argument (coordinate \(i\) of entry \(n\) freezes from stage \(n+i+1\) on) produces diagonal
  limit points \(\alpha_n\in\N\to\N\) such that every truncated family
  \(\bigl(f_n[\mathrm{cylinder}(\alpha_n,k-n)]\bigr)\) is non-separated, hence in
  particular nonempty. If all \(f_n(\alpha_n)\) coincided, that common value would lie in
  \(\bigcap_n\mathrm{range}\,f_n=\emptyset\); so two differ, and Hausdorff-separate them by open
  \(u\ni f_i(\alpha_i)\), \(v\ni f_j(\alpha_j)\). At a late enough stage the ultrametric cylinders
  around \(\alpha_i,\alpha_j\) shrink inside \(f_i^{-1}u\), \(f_j^{-1}v\), so
  \((u,v,\mathrm{univ},\dots)\) is a Borel separator of that truncated family --- contradicting its
  non-separatedness.
\end{proof}

\begin{theorem}[Generalized first separation theorem, Novikov; Srivastava 4.6.1]\label{thm:dst-first-separation}
  \lean{ErgodicTheory.generalized_first_separation}\leanok
  \uses{thm:dst-sepfam-ranges}
  A countable family \((A_n)\) of analytic sets with empty intersection is Borel separated: there
  are measurable \(B_n\supseteq A_n\) with \(\bigcap_n B_n=\emptyset\).
\end{theorem}
\begin{proof}\leanok
  \uses{thm:dst-sepfam-ranges}
  If some \(A_{n_0}=\emptyset\), take \(B_{n_0}=\emptyset\) and \(B_n=\mathrm{univ}\) otherwise.
  Otherwise every \(A_n\) is a nonempty continuous image of Baire space, \(A_n=\mathrm{range}\,f_n\),
  and \Cref{thm:dst-sepfam-ranges} applies verbatim.
\end{proof}

\section{The reduction ladder}

From the first separation theorem the ladder proceeds through the weak reduction principle to the
two section theorems, all on Polish spaces and their Borel structure --- no coanalytic pointclass
is developed.

\begin{theorem}[Weak reduction principle for coanalytic sets, Srivastava 4.6.5]\label{thm:dst-weak-reduction}
  \lean{ErgodicTheory.weak_reduction_coanalytic}\leanok
  \uses{thm:dst-first-separation, thm:dst-inter}
  If \((S_n)\) is a countable family with analytic complements (the \(S_n\) are coanalytic) and
  Borel union \(\bigcup_n S_n\), there are pairwise-disjoint measurable \(D_n\subseteq S_n\) with
  \(\bigcup_n D_n=\bigcup_n S_n\).
\end{theorem}
\begin{proof}\leanok
  \uses{thm:dst-first-separation, thm:dst-inter}
  Apply \Cref{thm:dst-first-separation} to the analytic sets \((\bigcup_m S_m)\cap(S_n)^{\mathsf c}\)
  (analytic by \Cref{thm:dst-inter}), whose intersection over \(n\) is empty. The resulting Borel
  separators \(D'_n\) give measurable pieces \((\bigcup_m S_m)\setminus D'_n\subseteq S_n\) with the
  correct union; disjointifying (\texttt{disjointed}) produces the pairwise-disjoint family.
\end{proof}

\begin{theorem}[Saint Raymond's theorem, Srivastava 4.7.1]\label{thm:dst-saint-raymond}
  \lean{ErgodicTheory.saintRaymond_closedSections}\leanok
  \uses{thm:dst-weak-reduction, thm:dst-inter, thm:dst-union, thm:dst-prod}
  For disjoint analytic \(A_0,A_1\subseteq X\times Y\) whose \(A_0\)-sections \(\{y\mid(x,y)\in
  A_0\}\) are closed, and a countable basis \((V_n)\) of \(Y\), there are Borel \(B_n\subseteq X\)
  with \(A_1\subseteq\bigcup_n B_n\times V_n\) and \(A_0\cap\bigcup_n B_n\times V_n=\emptyset\).
\end{theorem}
\begin{proof}\leanok
  \uses{thm:dst-weak-reduction, thm:dst-inter, thm:dst-union, thm:dst-prod}
  Replace \(A_1\) by a Borel Lusin separator \(A_1'\) from \(A_0\)
  (\texttt{AnalyticSet.measurablySeparable}). Let \(C_n=\{x\mid V_n\text{ misses the section
  }(A_0)_x\}\); its complement is the analytic projection \(\pi_X(A_0\cap(X\times V_n))\)
  (\Cref{thm:dst-inter}), and the closed sections give the basis decomposition
  \(A_0^{\mathsf c}=\bigcup_n C_n\times V_n\). Feed the coanalytic family
  \(S_n=A_1'\cap(C_n\times V_n)\) --- whose complements are analytic by \Cref{thm:dst-union} and
  \Cref{thm:dst-prod} and whose union is the Borel \(A_1'\) --- to the weak reduction principle
  (\Cref{thm:dst-weak-reduction}), then Lusin-separate each analytic projection \(\pi_X(D_n)\) from
  the analytic \((C_n)^{\mathsf c}\) to get the Borel \(B_n\).
\end{proof}

\begin{theorem}[Kunugui--Novikov theorem, Srivastava 4.7.2]\label{thm:dst-kunugui-novikov}
  \lean{ErgodicTheory.kunuguiNovikov_openSections}\leanok
  \uses{thm:dst-saint-raymond}
  A Borel \(B\subseteq X\times Y\) all of whose sections \(\{y\mid(x,y)\in B\}\) are open is a
  countable union of Borel rectangles \(B=\bigcup_n B_n\times V_n\) over any countable basis
  \((V_n)\) of \(Y\).
\end{theorem}
\begin{proof}\leanok
  \uses{thm:dst-saint-raymond}
  The special case \(A_0=B^{\mathsf c}\), \(A_1=B\) of Saint Raymond's theorem: closed sections of
  \(B^{\mathsf c}\) are exactly open sections of \(B\). Saint Raymond covers \(B\) by
  \(\bigcup_n B_n\times V_n\) while keeping it disjoint from \(B^{\mathsf c}\), forcing the reverse
  inclusion and hence equality.
\end{proof}

\section{The compact-section projection theorem}

The headline of the ladder is Novikov's theorem that a Borel set with compact sections has a Borel
projection (Srivastava 4.7.11; Kechris \S28, Arsenin--Kunugui). It is reached in two moves: the
compact-fibre-space case via a topology refinement making the set closed, then the general case by
compactifying the fibre through the Hilbert cube.

\begin{lemma}[Enumerated countable basis]\label{thm:dst-nat-basis}
  \lean{ErgodicTheory.exists_nat_basis}\leanok
  A second-countable space admits an \(\N\)-indexed basis \((V_n)\) (with possible empty entries):
  every point of every open set lies in some \(V_n\) contained in that open set.
\end{lemma}
\begin{proof}\leanok
  Insert \(\emptyset\) into a countable topological basis to make it a nonempty countable set, then
  enumerate it as a range.
\end{proof}

\begin{theorem}[Closed sections become closed after a topology refinement, Srivastava 4.7.4]\label{thm:dst-finer-polish}
  \lean{ErgodicTheory.exists_finer_polish_isClosed_of_closedSections}\leanok
  \uses{thm:dst-kunugui-novikov, thm:dst-nat-basis}
  For a Borel \(B\subseteq X\times Y\) with closed sections, there is a finer Polish topology
  \(t'\le t_X\) on \(X\) alone in which \(B\) becomes closed (in the product topology \(t'\times
  t_Y\)).
\end{theorem}
\begin{proof}\leanok
  \uses{thm:dst-kunugui-novikov, thm:dst-nat-basis}
  The complement \(B^{\mathsf c}\) has open sections, so Kunugui--Novikov
  (\Cref{thm:dst-kunugui-novikov}) writes it as \(\bigcup_n B_n\times V_n\) with the \(V_n\) from an
  enumerated basis (\Cref{thm:dst-nat-basis}). Each Borel \(B_n\) is \emph{clopenable}: there is a
  finer Polish topology making it clopen. Amalgamating these countably many refinements into a
  single Polish topology \(t'\) (\texttt{PolishSpace.exists\_polishSpace\_forall\_le}) makes every
  \(B_n\) open, hence \(B^{\mathsf c}=\bigcup_n B_n\times V_n\) open and \(B\) closed in
  \(t'\times t_Y\).
\end{proof}

\begin{theorem}[Compact-section projection, compact fibre space]\label{thm:dst-proj-compact}
  \lean{ErgodicTheory.measurableSet_image_fst_of_isCompact_sections_of_compactSpace}\leanok
  \uses{thm:dst-finer-polish}
  If \(Y\) is compact, a Borel \(B\subseteq X\times Y\) with compact sections has Borel projection
  \(\pi_X(B)\).
\end{theorem}
\begin{proof}\leanok
  \uses{thm:dst-finer-polish}
  Refine \(X\)'s topology to make \(B\) closed (\Cref{thm:dst-finer-polish}; compact sections are
  in particular closed). The first projection \(X\times Y\to X\) is a proper map when \(Y\) is
  compact, hence closed, so \(\pi_X(B)\) is closed in the refined topology \(t'\), thus Borel for
  \(t'\). Since \(t'\le t_X\) with both Polish, the two Borel \(\sigma\)-algebras coincide
  (\texttt{borel\_eq\_borel\_of\_le}, Lusin--Souslin), so \(\pi_X(B)\) is Borel for the original
  topology.
\end{proof}

\begin{lemma}[Continuous injection into the Hilbert cube]\label{thm:dst-cube}
  \lean{ErgodicTheory.exists_continuous_injective_toCube}\leanok
  Every Polish space \(Y\) admits a continuous injection into the compact Polish cube
  \(\N\to[0,1]\).
\end{lemma}
\begin{proof}\leanok
  For \(Y\) nonempty, fix a dense sequence \((u_n)\) in a compatible complete metric and send
  \(y\mapsto\bigl(\min(\mathrm{dist}(y,u_n),1)\bigr)_n\). This is continuous, and injective: if two
  points have equal coordinates then approximating each by the dense sequence forces their distance
  to \(0\). The empty case is vacuous.
\end{proof}

\begin{theorem}[Novikov compact-section projection theorem, Srivastava 4.7.11]\label{thm:dst-proj}
  \lean{ErgodicTheory.measurableSet_image_fst_of_isCompact_sections}\leanok
  \uses{thm:dst-proj-compact, thm:dst-cube}
  For Polish \(X,Y\), a Borel \(B\subseteq X\times Y\) all of whose sections are compact has Borel
  projection \(\pi_X(B)\) (Kechris, Thm.~28.8).
\end{theorem}
\begin{proof}\leanok
  \uses{thm:dst-proj-compact, thm:dst-cube}
  Compactify the fibre: push \(B\) into \(X\times(\N\to[0,1])\) along \(\mathrm{id}\times e\) for a
  continuous injection \(e\) (\Cref{thm:dst-cube}), which is a measurable embedding by
  Lusin--Souslin, so the image is Borel with the same first projection. The pushed sections are
  \(e[\,\text{compact section}\,]\), still compact hence closed in the compact cube, so the
  compact-fibre-space case (\Cref{thm:dst-proj-compact}) gives the Borel projection.
\end{proof}

\section{The everywhere-Borel singular Oseledets filtration}

We now discharge the payoff of issue \#11: the singular forward Lyapunov projector, previously
known only \texttt{AEMeasurable}, is \emph{everywhere} \texttt{Measurable}. The key geometric
observation is that closed-ball slices of a subspace fibre are compact, so the compact-section
projection theorem applies directly --- no \(\sigma\)-compact machinery is needed.

\begin{theorem}[Everywhere-Borel distance map via closed balls]\label{thm:dst-infdist}
  \lean{ErgodicTheory.measurable_infDist_of_measurableGraph}\leanok
  \uses{thm:dst-proj}
  Let \(V:X\to\mathrm{Submodule}\,\R\,(\R^d)\) over a standard Borel base \(X\) have a measurable
  graph \(\{(x,v)\mid v\in V x\}\). Then for every \(c\in\R^d\) the map \(x\mapsto\mathrm{infDist}\,
  c\,(V x)\) is \texttt{Measurable}.
\end{theorem}
\begin{proof}\leanok
  \uses{thm:dst-proj}
  By \texttt{measurable\_of\_Iic} it suffices that each sublevel \(\{x\mid\mathrm{infDist}\,c\,(V
  x)\le r\}\) is Borel. This set is the first projection of the graph sliced with the
  \emph{closed} ball, \(\{(x,v)\mid v\in V x\}\cap(X\times\overline B(c,r))\): a point \(x\) is in
  the sublevel iff the closed subspace \(V x\) --- which attains its distance to \(c\) --- meets
  \(\overline B(c,r)\). Each section \(V x\cap\overline B(c,r)\) is a closed subspace intersected
  with a compact ball in the proper Euclidean space, hence \emph{compact}, so the Novikov
  projection theorem (\Cref{thm:dst-proj}) makes the projection Borel.
\end{proof}

\begin{theorem}[Everywhere-measurable orthogonal projector from a measurable graph]\label{thm:dst-projmatrix-converter}
  \lean{ErgodicTheory.measurable_orthProjMatrix_of_measurableGraph}\leanok
  \uses{thm:dst-infdist, orthProjMatrix}
  Over a standard Borel base \(X\), a measurable subspace graph \(\{(x,v)\mid v\in V x\}\) makes the
  orthogonal-projection matrix \(x\mapsto\mathrm{orthProjMatrix}(V x)\) \texttt{Measurable}.
\end{theorem}
\begin{proof}\leanok
  \uses{thm:dst-infdist, orthProjMatrix}
  Reduce to entrywise measurability (\texttt{Matrix} carries the Pi structure). Each entry is a
  projected-basis coordinate (\texttt{orthProjMatrix\_apply}), which the polarisation identity
  \texttt{starProjection\_apply\_coord} writes as a fixed real combination of the three scalar
  distance maps \(x\mapsto\mathrm{infDist}\,c\,(V x)\) for \(c\) among the two basis vectors and
  their difference. Each is \texttt{Measurable} by \Cref{thm:dst-infdist}, and \texttt{Measurable}
  arithmetic with \texttt{measurable\_pi\_lambda} assembles the full projector.
\end{proof}

For contrast we record the a.e.\ tier, which the everywhere converter strictly strengthens.

\begin{theorem}[A.e.-measurable projector from a measurable graph]\label{thm:dst-aeconverter}
  \lean{ErgodicTheory.aemeasurable_orthProjMatrix_of_measurableGraph}\leanok
  \uses{thm:dst-nullmeas, orthProjMatrix}
  For any s-finite measure \(\mu\) on a standard Borel base, a measurable subspace graph makes
  \(x\mapsto\mathrm{orthProjMatrix}(V x)\) \texttt{AEMeasurable}.
\end{theorem}
\begin{proof}\leanok
  \uses{thm:dst-nullmeas, orthProjMatrix}
  The same polarisation reduction, but the distance sublevels are obtained only up to a null set,
  from the universal measurability of the analytic projection (\Cref{thm:dst-nullmeas}) rather than
  its honest Borel-ness. This is the route of Chapter~\ref{chap:lyapunov}; the everywhere converter
  (\Cref{thm:dst-projmatrix-converter}) replaces \Cref{thm:dst-nullmeas} by the Novikov projection
  theorem.
\end{proof}

\begin{theorem}[Issue \#11 headline: everywhere-Borel forward Lyapunov projector]\label{thm:dst-headline}
  \lean{ErgodicTheory.measurable_orthProjMatrix_lambdaSublevel}\leanok
  \uses{thm:dst-projmatrix-converter, lambdaSublevel, lambdaBar, IsUltrametricGrowth}
  Let \(A:X\to\Matrix_d(\R)\) be a measurable generator and \(T:X\to X\) measurable over a standard
  Borel base \(X\), and assume the \emph{everywhere} ultrametric-growth gate
  \(\forall x,\ \mathrm{IsUltrametricGrowth}(\mathrm{lambdaBar}\,A\,T\,x)\). Then for every
  threshold \(c\) the sublevel projector \(x\mapsto\mathrm{orthProjMatrix}(\mathrm{lambdaSublevel}\,
  A\,T\,x\,c)\) is \emph{everywhere} \texttt{Measurable}. No invertibility (\(\det\neq 0\)),
  ergodicity, measure-preservation, or log-norm integrability is assumed --- \(A\) and \(T\) are
  unrelated, the invertible-MET hypotheses traded for the everywhere gate.
\end{theorem}
\begin{proof}\leanok
  \uses{thm:dst-projmatrix-converter}
  The sublevel filtration has a measurable graph
  \lean{ErgodicTheory.measurableSet_graph_lambdaSublevel}
  built from the measurable data \(A,T\) and the everywhere gate, and the general converter
  (\Cref{thm:dst-projmatrix-converter}) turns that graph into the everywhere-measurable projector
  via the Novikov compact-section projection theorem.
\end{proof}

\begin{theorem}[The a.e.\ sublevel projector, issue \#6]\label{thm:dst-aemeasurable}
  \lean{ErgodicTheory.aemeasurable_orthProjMatrix_lambdaSublevel}\leanok
  \uses{thm:dst-aeconverter, lambdaSublevel, isUltrametricGrowth_lambdaBar}
  For the invertible-MET data (ergodic, measure-preserving \(T\), integrable log-norms), the
  sublevel projector \(x\mapsto\mathrm{orthProjMatrix}(\mathrm{lambdaSublevel}\,A\,T\,x\,c)\) is
  \texttt{AEMeasurable}.
\end{theorem}
\begin{proof}\leanok
  \uses{thm:dst-aeconverter, isUltrametricGrowth_lambdaBar}
  Here the ultrametric-growth gate holds only a.e.\ (\texttt{isUltrametricGrowth\_lambdaBar},
  itself needing invertibility, ergodicity, and integrable log-norms), so the a.e.\ converter
  (\Cref{thm:dst-aeconverter}) applies. The everywhere headline
  (\Cref{thm:dst-headline}) is strictly stronger: it drops every dynamical hypothesis in exchange
  for the pointwise gate.
\end{proof}

\section{Descriptive complexity of the continuous-section seal (issue \#61)}

The projection machinery above answers a \emph{measurability} question; the same analytic /
coanalytic pointclass vocabulary answers a \emph{complexity} question about ergodic-theoretic
invariants. Fix a compact metric (Borel) carrier \(X\) and the Polish parameter space
\[
  \mathrm{Params}\,X \;=\; C(X,X)^3 \times P(X)^2,
\]
whose points are tuples \((T,S,\pi,\mu,\nu)\): a base map, a factor map, a factor projection, and the
two probability measures. A candidate \emph{continuous} section \(s \in C(X,X)\) is admissible when
\(\pi\circ T = S\circ\pi\), \(\pi_*\mu = \nu\), \(\pi\circ s = \mathrm{id}\), \(s\circ S = T\circ s\),
and \(s_*\nu = \mu\). Whether a system \emph{admits} such a section is the continuous,
compact-metric-parametrised reading of the Foreman--Rudolph--Weiss section-existence problem; we
place it in the projective hierarchy over the Polish space \(\mathrm{Params}\,X\).

Two Polish-space infrastructure facts, both genuine gaps in Mathlib's weak-convergence API, make the
hierarchy apply and are natural upstream candidates.

\begin{theorem}[The space of probability measures on a compact metric space is Polish]\label{thm:dst-pspace-polish}
  \lean{MeasureTheory.polishSpace_probabilityMeasure}\leanok
  For a compact metric Borel space \(X\), the space \(P(X) = \mathrm{ProbabilityMeasure}\,X\) with the
  topology of convergence in distribution is Polish (Kechris \S17.E).
\end{theorem}
\begin{proof}\leanok
  Prokhorov's theorem makes \(P(X)\) compact and the L\'evy--Prokhorov metric makes it metrizable and
  second countable. A compact metrizable space is completely metrizable
  (\texttt{isCompletelyMetrizableSpace\_of\_compactSpace}: a compatible metric on a compact --- hence
  complete --- uniform space is complete), which upgrades the package to Polishness. A companion
  \texttt{polishSpace\_prod} assembles the finite products.
\end{proof}

\begin{theorem}[Joint continuity of the pushforward]\label{thm:dst-pushforward-cont}
  \lean{MeasureTheory.continuous_probabilityMeasure_map_compact}\leanok
  For a compact metric \(X\) and a metric \(Y\), the pushforward
  \((f,\nu) \mapsto f_*\nu : C(X,Y) \times P(X) \to P(Y)\) is jointly continuous, with
  \(C(X,Y)\) carrying the uniform metric and \(P(X)\) the topology of convergence in distribution.
\end{theorem}
\begin{proof}\leanok
  The Billingsley mapping-theorem argument. Test against a bounded continuous \(g : Y \to \R\) via
  \(\int g\,\mathrm d(f_*\nu) = \int (g\circ f)\,\mathrm d\nu\) and split the difference into a
  weak-convergence term (\(\nu_i \to \nu\) tested against \(g\circ\sigma\)) and a uniform term bounded
  by \(\|g\circ\sigma_i - g\circ\sigma\|_\infty\), which vanishes because postcomposition
  \(\tau \mapsto g\circ\tau : C(X,Y) \to (X \to^b \R)\) is continuous. The filter form is
  \texttt{tendsto\_probabilityMeasure\_map\_of\_tendsto}. Mathlib previously had only the fixed-map
  continuity \texttt{ProbabilityMeasure.continuous\_map}.
\end{proof}

\begin{definition}[The section relation]\label{thm:dst-sectionrel}
  \lean{ErgodicTheory.SectionRel}\leanok
  For \(p = (T,S,\pi,\mu,\nu) \in \mathrm{Params}\,X\) and \(s \in C(X,X)\), the predicate
  \(\mathrm{SectionRel}\,p\,s\) is the conjunction of the five equalizer conditions
  \(\pi\circ T = S\circ\pi\), \(\pi_*\mu = \nu\), \(\pi\circ s = \mathrm{id}\), \(s\circ S = T\circ
  s\), \(s_*\nu = \mu\) making \(s\) a continuous, measure-preserving, equivariant section of \(\pi\).
  A parameter is \emph{sealed} (\texttt{ErgodicTheory.IsSealed}) when no such \(s\) exists.
\end{definition}

\begin{theorem}[The section relation is closed]\label{thm:dst-sectionrel-closed}
  \lean{ErgodicTheory.isClosed_sectionRel}\leanok
  \uses{thm:dst-sectionrel, thm:dst-pushforward-cont}
  \(\{(p,s) \mid \mathrm{SectionRel}\,p\,s\}\) is a closed subset of
  \(\mathrm{Params}\,X \times C(X,X)\).
\end{theorem}
\begin{proof}\leanok
  \uses{thm:dst-pushforward-cont}
  Each of the five conjuncts is an equalizer of jointly continuous maps into a Hausdorff space:
  composition is jointly continuous (\texttt{ContinuousMap.continuous\_comp'}), and the pushforward
  is jointly continuous (\Cref{thm:dst-pushforward-cont}); a finite intersection of closed sets is
  closed.
\end{proof}

\begin{theorem}[Section-existence is analytic; sealedness coanalytic]\label{thm:dst-section-analytic}
  \lean{ErgodicTheory.sectionExists_analyticSet}\leanok
  \uses{thm:dst-sectionrel-closed, thm:dst-pspace-polish}
  The set \(\{p \mid \exists s,\ \mathrm{SectionRel}\,p\,s\}\) of parameters admitting some continuous
  section is \emph{analytic} (\(\Sigma^1_1\)): it is the continuous image, under the coordinate
  projection \(\mathrm{Prod.fst}\), of the closed relation of \Cref{thm:dst-sectionrel-closed} inside
  the Polish space \(\mathrm{Params}\,X\). Dually the sealed set \(\{p \mid \mathrm{IsSealed}\,p\}\) is
  \emph{coanalytic} (\(\Pi^1_1\), \texttt{ErgodicTheory.isSealed\_coanalyticSet}). The identity
  parameter certifies non-vacuity (\texttt{ErgodicTheory.sectionExists\_nonempty}).
\end{theorem}
\begin{proof}\leanok
  \uses{thm:dst-sectionrel-closed, thm:dst-pspace-polish}
  A closed subset of a Polish space is analytic and analyticity is preserved by continuous images
  (here the coordinate projection \(\mathrm{Prod.fst}\)), giving the \(\Sigma^1_1\) bound; that
  \(\mathrm{Params}\,X\) and \(C(X,X)\) are Polish is \Cref{thm:dst-pspace-polish} and the uniform
  metric on \(C(X,X)\). The sealed set is by definition the complement, hence \(\Pi^1_1\); the
  identity tuple satisfies \(\mathrm{SectionRel}\) with \(s = \mathrm{id}\).
\end{proof}

\begin{remark}[Honest scope]\label{thm:dst-section-scope}
  This is the issue's sanctioned ``restricted class first'' reading: \emph{continuous} sections over
  \emph{compact-metric-parametrised} systems, for which \(C(X,X)\), \(P(X)\), and their products are
  all Polish, so the descriptive hierarchy applies verbatim. Two frontiers are disclosed, not
  delivered. The \emph{classical} Foreman--Rudolph--Weiss section-existence problem (Ann.\ of Math.\
  \textbf{173}, 2011) is stated for \emph{arbitrary measurable} sections over the \(L^0\)/MALG
  parametrisation, which needs \(L^0\) (measurable maps mod null sets) as a Polish space ---
  infrastructure Mathlib lacks. The \emph{tier-3} hardness statement (that the sealed set is
  \(\Sigma^1_1\)-complete, hence non-Borel) needs the Borel-reduction and \(\Sigma^1_1\)-completeness
  machinery of Kechris \S14, \S27, also absent from Mathlib. A concrete sealed witness is left to its
  own follow-up.
\end{remark}

\section{An honest frontier: the general \texorpdfstring{\(\sigma\)}{sigma}-compact-section theorem}

The compact-section projection theorem (\Cref{thm:dst-proj}) suffices for the subspace-valued
Oseledets filtration precisely because closed-ball slices of the subspace fibres are compact
(\Cref{thm:dst-infdist}). The fully general theorem --- that a Borel set with
\(\sigma\)-\emph{compact} sections has Borel projection (the Arsenin--Kunugui theorem, Saint
Raymond 5.12.2; Kechris \S28) --- is a strictly deeper result whose formalization would require the
Effros Borel structure on the hyperspace of closed sets, \(\Pi^1_1\)-boundedness, and the
transfinite rank analysis of coanalytic sets. It is \emph{not} formalized in this development, and
is recorded here as a documented research frontier. The everywhere-Borel singular filtration of
this chapter deliberately routes around it: the closed-ball compactification of the fibre keeps the
whole issue \#11 payoff inside the compact-section case, so no \(\sigma\)-compact machinery is ever
invoked.
